\documentclass{anvil-paper}
\usepackage{array}
\usepackage{tabularx}
\usepackage{fontspec}
\setmainfont{Times New Roman}
\setmonofont{Menlo}[Scale=0.85]
\usepackage[htt]{hyphenat}

% Lean identifiers: typewriter with url-style breaking at underscores and dots,
% plus a discouraged (penalty 500) break before each capital letter so that long
% camel-case identifiers can break without a badness-10000 line.
\makeatletter
\def\lean@camel{%
  \do\A{\penalty500\mathchar`A}\do\B{\penalty500\mathchar`B}\do\C{\penalty500\mathchar`C}%
  \do\D{\penalty500\mathchar`D}\do\E{\penalty500\mathchar`E}\do\F{\penalty500\mathchar`F}%
  \do\G{\penalty500\mathchar`G}\do\H{\penalty500\mathchar`H}\do\I{\penalty500\mathchar`I}%
  \do\J{\penalty500\mathchar`J}\do\K{\penalty500\mathchar`K}\do\L{\penalty500\mathchar`L}%
  \do\M{\penalty500\mathchar`M}\do\N{\penalty500\mathchar`N}\do\O{\penalty500\mathchar`O}%
  \do\P{\penalty500\mathchar`P}\do\Q{\penalty500\mathchar`Q}\do\R{\penalty500\mathchar`R}%
  \do\S{\penalty500\mathchar`S}\do\T{\penalty500\mathchar`T}\do\U{\penalty500\mathchar`U}%
  \do\V{\penalty500\mathchar`V}\do\W{\penalty500\mathchar`W}\do\X{\penalty500\mathchar`X}%
  \do\Y{\penalty500\mathchar`Y}\do\Z{\penalty500\mathchar`Z}}
\DeclareUrlCommand\lean{\urlstyle{tt}%
  \expandafter\def\expandafter\UrlSpecials\expandafter{\UrlSpecials\lean@camel}}
% Table variant: breaks only at underscores and dots.
\DeclareUrlCommand\leantab{\urlstyle{tt}}
\makeatother
% Names from Lean's standard library (not declarations of this project).
\newcommand{\stdlean}[1]{\lean{#1}}
% Public links to the branch of record (BRIEF rule R-LINK). \repobase is the one
% base URL, switched from the branch to the release tag at publication.
\newcommand{\repobase}{https://github.com/rjwalters/lean-genius/blob/erdos85/integration}
\newcommand{\repofile}[2]{\href{\repobase/#1}{\lean{#2}}}
\newcommand{\repofiletab}[2]{\href{\repobase/#1}{\leantab{#2}}}
\newcommand{\rp}{research/problems/erdos-85-wip-01}
\tolerance=800
\setlength{\emergencystretch}{3em}
\renewcommand{\topfraction}{.9}
\renewcommand{\bottomfraction}{.9}
\renewcommand{\textfraction}{.05}
\renewcommand{\floatpagefraction}{.85}
\newcommand{\AREG}{\mathrm{A\text{-}REG}}
\newcommand{\NONBIP}{\mathrm{A\text{-}REG\text{-}NONBIP}}
\newcommand{\Qn}{\textsf{Erdos85Question}}
\newcommand{\hnoFortyNine}{\lean{hno49}}
% Tool names (not Lean declarations).
\newcommand{\cakelpr}{\texttt{cake\textunderscore lpr}}
\newcommand{\vtwocnf}{\texttt{v2cnf}}

\newtheorem*{resultA}{Result A}
\newtheorem*{theoremB}{Theorem B}
\theoremstyle{definition}
\newtheorem*{hypAREG}{Hypothesis A-REG}
\newtheorem*{nonbipconn}{NONBIP-CONNECTED}

\title{A Drop at Forty-Nine in Erdős Problem 85}
\author{Claude Fable, GPT Sol, Astra, Claude Opus and Robb Walters \\ Lean Genius project, 2AM Logic}
\date{7 October 2026}

\begin{document}
\maketitle

\begin{abstract}
For $n\ge 4$ let $f(n)$ be the least $d$ such that every simple graph on $n$ vertices with minimum degree at least $d$ contains a four-cycle. Erdős Problem 85 asks whether $f$ is eventually nondecreasing. We report two results.

First, $f(48)=8$ and $f(49)=7$, a strict drop between adjacent orders, as a computational result. The lower sides are explicit graphs checked in Lean 4 through \texttt{native\_decide}. For the upper side, Lean proves that a $C_4$-free graph on 49 vertices with minimum degree 7 lies in one of four strata and, using \texttt{native\_decide} for a finite enumeration, reduces the largest, H1, to the unsatisfiability of 13,351 symmetry-reduced formulas. H1 is certificate-checked: the CakeML-verified checker \cakelpr{} accepted an unsatisfiability proof for each, about 28 TB of LRAT in all; the hardest was split into 36 cubes whose composition is a standard-axiom Lean lemma. Among the smaller strata, only one part has certificates, checked inside Lean (H7, $t\ge 1$); the others are closed by reviewed arguments with independently checked computation (H3, H5; H7, $t=0$), on which the upper bound also depends. The result is not a Lean theorem: for H1 the external checks are not admitted into Lean and each checked file is tied to its Lean formula only by a compiled Lean emitter; the H3, H5 and H7 ($t=0$) exclusions are open in Lean. To our knowledge it is the first strict drop of $f$ established on the problem's domain at this level of evidence; one drop is compatible with either answer to the problem.

Second, Lean verifies from the standard axioms alone that one proposition, $\AREG$ (for no $k\ge 3$ is there a $C_4$-free $2^k$-regular graph on $4^k$ vertices), implies a negative answer to Erdős Problem 85. $\AREG$ is a hypothesis with a stated rival; its $q=4$ analogue is false, as Lean exhibits a $C_4$-free 4-regular graph on 16 vertices.
\end{abstract}

\section{Introduction}\label{sec:intro}

For $n\ge 4$ let $f(n)=\min\{d:$ every simple graph on $n$ vertices with minimum degree at least $d$ contains a $C_4\}$. Erdős Problem 85 \citep{bloom-erdos85} asks whether $f(n+1)\ge f(n)$ for all sufficiently large $n$. In Lean 4 \citep{moura2021lean4}, on Mathlib's simple graphs \citep{mathlib2020}, the function is \lean{minDegreeForC4} and the question is the proposition \Qn; Lean proves that its negation is equivalent to the existence of arbitrarily late strict drops (\lean{erdos85Negation_iff_not_question}). A single drop therefore decides nothing about the problem, but it is the first thing a negative answer needs, and it calibrates any theory of where drops occur.

The problem has a Ramsey form. With Boza's convention $r(s)=R(C_4,K_{1,s})$ \citep{boza2024ramsey,zhang2017polarity}, $r(s)\le N$ holds exactly when $f(N)\le N-s$, so a plateau $r(s)=r(s+1)=N$ is the same as a drop $f(N)<f(N-1)$. Boza's table gives $r(41)=49$ and leaves $r(42)\in\{49,50\}$ open; our upper side selects $r(42)=49$, which is exactly $f(49)=7<8=f(48)$. We use $f$ throughout.

\begin{resultA}
$f(48)=8$ and $f(49)=7$; in particular $f(49)<f(48)$. The lower sides are Lean-checked. The upper side, $f(49)\le 7$, is the Lean hypothesis \hnoFortyNine\ $:\ \neg\,$\texttt{C4FreeMinDegreeWitness 49 7}, which we support stratum by stratum: by certificate-checked computation for H1 and for H7 with $t\ge 1$, and by reviewed arguments with independently checked computation for H3, H5 and H7 with $t=0$ (Table~\ref{tab:status}).
\end{resultA}

\begin{theoremB}[Lean, standard axioms]
Let $\AREG$ (\lean{BinarySquareRegularExclusion}) state that for every $k\ge 3$ there is no simple $C_4$-free $2^k$-regular graph on $4^k$ vertices. Then $\AREG\Rightarrow\neg\,\Qn$ (\lean{not_erdos85Question_of_binarySquareRegularExclusion}).
\end{theoremB}

\paragraph{What is new.} (1) The drop. To our knowledge no strict drop of $f$ on the problem's domain $n\ge 4$ was previously settled: the maintained record lists no partial result and warns that its literature list may be incomplete \citep{bloom-erdos85}, and Boza's table reports no consecutive equality that yields one, though some of its entries are still ranges \citep{boza2024ramsey}. (2) The check of H1. All 13,351 instances of the largest stratum carry unsatisfiability proofs accepted by the formally verified checker \cakelpr{}, about 28 TB in all, produced and checked without storing the fresh proofs; the one instance that no whole-instance solve finished is closed by a cube partition whose composition is a Lean lemma. Section~\ref{sec:h1} describes the method, positions it against its closest precedent, the formally verified empty-hexagon theorem \citep{heule2024hexagon,subercaseaux2024hexagonlean}, and explains how a reader can repeat the check. (3) Theorem B, which turns the infinite problem into one proposition, together with what is proved beneath it and a map of approaches that fail.

\paragraph{What is not claimed.} Result A is not a theorem of Lean: the H1 checks run outside the Lean kernel and are not admitted into Lean, the checked files are tied to the Lean formulas by a compiled emitter rather than a kernel proof, the exclusions of H3, H5 and H7 ($t=0$) rest on reviewed arguments rather than on certificates or Lean proofs, and Table~\ref{tab:status} lists, once, every step that separates Result A from a Lean theorem. Nothing is claimed about Erdős Problem 85 itself. $\AREG$ is not asserted; Section~\ref{sec:theoremB} gives the evidence on both sides. Every file and Lean module named here is linked to its public location (Section~\ref{sec:availability}), and every number traces to a linked receipt.

\section{The drop at forty-nine}\label{sec:drop}

\subsection{Lower sides}\label{sec:lower}

The module \repofile{proofs/Proofs/Erdos85FiniteDropWitnesses.lean}{Erdos85FiniteDropWitnesses.lean} contains a $C_4$-free 7-regular graph on 48 vertices (168 edges) and a $C_4$-free graph on 49 vertices with minimum degree 6. From them Lean proves $f(48)=8$ (\lean{minDegreeForC4_fortyEight_eq_eight_checked}) and $f(49)\ge 7$ (\lean{seven_le_minDegreeForC4_fortyNine_checked}), and, taking \hnoFortyNine\ as a hypothesis, $f(48)=8\wedge f(49)=7$ and $f(49)<f(48)$ (\lean{minDegreeForC4_fortyEight_fortyNine_exact_checked}, \lean{minDegreeForC4_fortyNine_lt_fortyEight_checked}). The graphs are checked through \texttt{native\_decide}. A cold rebuild from source (pinned Lean 4.31.0 image, fresh build volume) printed, for each witness, three per-declaration axioms of the form \texttt{\ldots\_native.native\_decide.ax\_N} beyond \texttt{propext}, \texttt{Classical.choice} and \texttt{Quot.sound}; these express trust in the compiled evaluator. The audit output prints only these per-declaration names; identifying them with \texttt{Lean.ofReduceBool} is a statement about how Lean implements \texttt{native\_decide}, not a line of that output. Theorem B, in the same audit, uses the three standard axioms only. The 48-vertex witness is not isomorphic to any of the ten graphs with 48 vertices and 168 edges in the Afzaly--McKay records \citep{afzaly-mckay-extremal} (\repofile{\rp/sat49/verify_boza48_nonisomorphism.py}{verify_boza48_nonisomorphism.py}); those records also contain 49-vertex graphs of minimum degree 6, which corroborate $f(49)\ge 7$ but say nothing about the upper side.

\subsection{The case split}\label{sec:split}

Let $G$ be $C_4$-free on 49 vertices with minimum degree at least 7. Two vertices share at most one neighbour, so the non-returning walks of length two from $v$ end at distinct vertices and $\sum_{u\sim v}(\deg u-1)\le 48$. Hence every degree is 7 or 8, and a vertex of degree 8 (\emph{high}) has only neighbours of degree 7. The degree sum $7\cdot 49+h$ is even, so the number $h$ of high vertices is odd, and Lean proves $h\in\{1,3,5,7,9\}$ (\lean{orderFortyNine_card_high_eq_one_or_three_or_five_or_seven_or_nine}). The case $h=9$ is refuted in Lean with no external input (\lean{orderFortyNineStratumExcluded_nine}), and
\[
\begin{aligned}
&\text{\lean{not_c4FreeMinDegreeWitness_fortyNine_seven_of_strata}}\ (h_1)\,(h_3)\,(h_5)\,(h_7)\\
&\qquad:\ \neg\,\texttt{C4FreeMinDegreeWitness 49 7}
\end{aligned}
\]
takes exactly the four exclusions of the strata H1, H3, H5 and H7. The case split is therefore part of the formal statement, and the evidence below is evidence for those four propositions.

\subsection{The three small strata}\label{sec:small}

Write the high vertices first, so that the adjacency matrix is $\left[\begin{smallmatrix}0&B\\B^{\mathsf T}&C\end{smallmatrix}\right]$. High vertices are pairwise non-adjacent, any two of them have exactly one common neighbour, and a low vertex has at most three high neighbours, because its neighbourhood induces a matching. Up to relabelling, the high--low incidence falls into finitely many \emph{cells}, indexed by the number $t$ of low vertices with three high neighbours.

\paragraph{H3.} Any two of the three high vertices have exactly one common neighbour, so there are two cells: in $t=1$ one low vertex is adjacent to all three high vertices, and in $t=0$ each pair of high vertices has its own common neighbour, these three low vertices spanning a matching with $b\in\{0,1\}$ edges. Lean proves that every graph of the stratum, after relabelling, satisfies the constraints of one of the two cell representatives (\lean{threeHighCanonicalGraphCover_all} in \repofile{proofs/Proofs/Erdos85OrderFortyNineThreeHighOneFiber.lean}{Erdos85OrderFortyNineThreeHighOneFiber.lean}), and \lean{orderFortyNineStratumExcluded_three_of_representativeExclusions} reduces the stratum to the exclusion of the two representatives. That exclusion is not certificate-checked: no SAT verdict or LRAT proof exists for either cell formula, and in Lean the two exclusions are undischarged hypotheses. H3 is excluded, outside Lean, by a reviewed paper argument (\repofile{\rp/Q7_H3_PROFILE_EXCLUSION_20260910.md}{Q7_H3_PROFILE_EXCLUSION_20260910.md}): graph-to-core reductions followed by exhaustive enumerations in Python over 3,337 induced cases for $t=1$, and over 972 ($b=0$) and 3,600 ($b=1$) combinations of normalized core and host choice for $t=0$; every case is excluded. Each enumeration passed independent review, and the two $t=0$ searches were replayed from unchanged source with no time limit, in exact agreement with the retained case records.

\paragraph{H5.} There are three cells, $t=0,1,2$. They were closed by reviewed reductions with independently checked finite computation, joined by a reviewed ledger (\repofile{\rp/q7_h5_closure_ledger/README.md}{q7_h5_closure_ledger}) whose checker reconstructs all 49 masks of each cell against the Lean canonical arrays. In the hardest cell, T2, 13 heavy classes remain after the reductions; 12 are excluded by a reviewed census chain, and the remaining core is split into 92 branches, each assigned to a reviewed exclusion. No SAT certificate is involved. The Lean interface \lean{orderFortyNineStratumExcluded_five_of_booleanExclusions} takes three Boolean-exclusion premises, which are not yet discharged in Lean.

\paragraph{H7.} The seven high vertices carry a linear triple system with $t\le 7$ triples (\lean{orderFortyNine_highIncidence_profile_of_seven_high}). A kernel-checked census shows that every such system is equivalent to one of fourteen canonical representatives, and Lean proves that every graph of the stratum satisfies the constraints of one of them (\lean{sevenHighCanonicalGraphCover_all}). The thirteen representatives with $t\ge 1$ are excluded by LRAT certificates checked inside Lean by \texttt{native\_decide} against the formula \lean{orderFortyNineGeneratedH7SatCnf} (1,329,041 clauses each), and \lean{orderFortyNineStratumExcluded_seven_of_t0} reduces the stratum to the representative with $t=0$. These thirteen proofs depend on \texttt{Lean.ofReduceBool}; they were compiled when the certificates were produced but not rebuilt in the cold audit, and the packed certificates they read are not published.

In the $t=0$ cell the 42 low vertices split into 7 with no high neighbour, 14 with one and 21 with two, and the argument works from the graph induced on the seven: it is $C_4$-free with maximum degree 3 and has 43 isomorphism classes with 6 to 9 edges. A counting argument on common-neighbour pairs excludes 15 classes; a reviewed closure ledger (\repofile{\rp/H7_CLOSURE_20260915.md}{H7_CLOSURE_20260915.md}) excludes the other 28 by structural arguments supported by finite computation; the completeness of one source enumeration rests on an audit of the enumerator code rather than an independent replay. The largest of these checks partitions 2,278,608 leaves, and its final 445,699 certificates were reproduced by a third, independently written implementation with no mismatch. The supporting capacity inequalities are Lean theorems; the closure itself is not, and its Lean capstone \lean{orderFortyNineStratumExcluded_seven_of_emptyCubeEvidenceVectors} is not yet instantiated.

\subsection{H1 and its 13,351 orbits}\label{sec:h1-structure}

H1, one high vertex, carries most of the computation. The eight neighbours of the high vertex induce a perfect matching, every other vertex has exactly one neighbour among them, and the remaining 40 vertices induce a 6-regular $C_4$-free graph with eight attachment groups of five. Such a graph determines a profile (one of five) and a \emph{miss table} of 24 small counts that record, for pairs of the high vertex's neighbours, how the attachment groups meet. Permutations of the eight neighbours that preserve the matching and the profile act on tables, and Lean proves that every H1 graph can be labelled so that its table is the stored representative of its orbit; so one formula per orbit suffices. An \emph{orbit} below is such a representative together with its SAT formula, the Lean term \lean{oneHighFamilyV2SatCnf} at that table (the hardest has 42,160 variables and 613,228 clauses).

The stored inventory, which Lean reads as a file, has 13,541 orbit representatives; a capacity inequality that Lean proves for every H1 graph removes 190 of them, leaving 13,351 (\lean{oneHighCapacityInventory_total_length}, by \texttt{native\_decide}). Lean then proves
\[
\begin{aligned}
&\bigl(\forall\ \text{orbits } (p,T) \text{ of the 13,351}:\ \texttt{OneHighFamilyV2CheckedUnsat}\ p\ T\bigr)\\
&\qquad\Longrightarrow\ \texttt{OrderFortyNineStratumExcluded 1}
\end{aligned}
\]
as \lean{orderFortyNineStratumExcluded_one_of_capacityInventory_checked}, where \lean{OneHighFamilyV2CheckedUnsat} says that the orbit's Lean formula is unsatisfiable. Its cover argument (every H1 graph induces, after relabelling, one of the stored tables) uses \texttt{native\_decide} for the finite enumeration checks: its printed axioms are \texttt{propext}, \texttt{Classical.choice} and \texttt{Quot.sound} together with 23 \texttt{native\_decide} axioms from those checks, and no \texttt{sorryAx} (\repofile{\rp/h1_cert_pilot_20261001/H1_COVER_AXIOMS_20261006.txt}{H1_COVER_AXIOMS_20261006.txt}). So, in Lean, H1 reduces to 13,351 SAT questions, and Section~\ref{sec:h1} answers all of them.

Before any proof was logged, the 1,161 orbits that had no certificate from the first production run in August 2026 were settled verdict-only by a two-solver census (Kissat 4.0.4 \citep{biere2024kissat}, then CaDiCaL 3.0.1 \citep{biere2020cadical,biere2024cadical2}; about 5,700 core-hours, no satisfiable instance, no disagreement; \repofile{\rp/phase_b_h1_census_20260927/CENSUS.md}{CENSUS.md}). That census is now superseded by certificates and we do not rely on it.

\section{How H1 was checked}\label{sec:h1}

A solver's UNSAT verdict is not a proof, but it can be made into one. DRAT proofs are checked by drat-trim \citep{wetzler2014drat}; the LRAT format adds hints that let a small checker replay each step \citep{cruzfilipe2017lrat}; and \citet{tan2021cakelpr} give \cakelpr{}, an LRAT checker whose correctness, including file parsing, is proved in HOL4 and carried to machine code by the verified CakeML compiler. The largest SAT-based results in combinatorics check proofs too large to keep \citep{heule2016pythagorean,heule2018schur}.

The closest precedent for what we do is the empty-hexagon theorem. \citet{heule2024hexagon} showed by cube-and-conquer \citep{heule2011cube} that every set of 30 points in general position in the plane contains an empty convex hexagon, with the proof checked by \cakelpr{}, and \citet{subercaseaux2024hexagonlean} verified in Lean that the SAT encoding is faithful to the geometric statement and joined that verification to the \cakelpr{} check. The hexagon verification and ours both reduce the mathematical statement, in Lean, to the unsatisfiability of a formula defined in Lean, and both rely on \cakelpr{} for the unsatisfiability check. On this axis we claim no advantage: our reduction uses 23 \texttt{native\_decide} axioms (Section~\ref{sec:h1-structure}), and we tie each checked CNF file to its Lean formula only through a compiled Lean emitter (Section~\ref{sec:trust}). What this paper adds is of a different kind: a Lean reduction of a whole stratum to 13,351 formulas, all checked by \cakelpr{} (about 28 TB of proofs); a check-then-discard discipline that keeps only hashes and relies on byte-identical regeneration; the re-validation of an archived certificate bank of 22.55 TB; and a public checker with which a reader can re-check any H1 orbit.

\subsection{Check-then-discard}\label{sec:method}

For each orbit the pipeline does three things.
\begin{enumerate}
\item \emph{Pin the formula.} Regenerate the CNF from the orbit's table with \vtwocnf, the compiled form of the Lean emitter \repofile{proofs/Proofs/Erdos85OneHighV2CnfEmit.lean}{Erdos85OneHighV2CnfEmit.lean}, which prints the Lean term \lean{oneHighFamilyV2SatCnf} as DIMACS; run its \texttt{check} mode; and require the CNF's sha256 to equal the hash recorded when the orbit was first solved. The emitter binary and container image are pinned by hash.
\item \emph{Check the proof as it streams.} Pipe an LRAT proof, fresh from the solver or read from an archive, through a hashing relay into \cakelpr{} (commit \texttt{a36874a8}). An orbit counts only when \cakelpr{} prints \texttt{s VERIFIED UNSAT} and, for a fresh solve, the solver reports \texttt{s UNSATISFIABLE}. \cakelpr{} exits with status 0 even when a check fails, so that line is the only success signal.
\item \emph{Discard.} Keep the proof's sha256 and length, not the proof. With the identical solver binary the proofs are reproducible byte for byte: three orbits solved twice on different machine types gave identical proofs of 27 to 28.6 GB. A macOS build of the same solver release gives different, equally valid proofs, so reproduction pins the exact binary.
\end{enumerate}
\cakelpr{} reads the proof with a heap fixed in advance, whatever the proof's length: 4 GB sufficed for 1,075 of 1,160 fresh proofs and none needed more than 16 GB, and checking cost about 6\% of the solver's CPU time. This bounded memory is what made whole-instance certificates practical; before the run we had planned to cube every hard orbit to keep proofs small.

\subsection{Four evidence sets}\label{sec:sets}

The 13,351 orbits fall into four disjoint sets (Table~\ref{tab:h1}); their union is the capacity inventory, with no orbit missing.

\begin{table}[htbp]
\centering\small
\caption{The H1 check. Every orbit has an unsatisfiability proof accepted by the CakeML-verified checker \cakelpr{}; no proof was rejected and no orbit was satisfiable. All checks ran outside the Lean kernel.}
\label{tab:h1}
\begin{tabularx}{\textwidth}{@{}>{\raggedright\arraybackslash}p{.20\textwidth}>{\raggedleft\arraybackslash}p{.07\textwidth}>{\raggedright\arraybackslash}X>{\raggedright\arraybackslash}p{.24\textwidth}@{}}
\toprule
Set & Orbits & Proofs & Checker \\
\midrule
August 2026 bank & 12,094 & archived Kissat proofs, converted to LRAT by drat-trim; 22.55 TB, largest 26.5 GB & \cakelpr{}, 8 GB heap \\
Fresh census & 1,160 & CaDiCaL 3.0.1 binary LRAT, streamed; 5.64 TB, largest 36.2 GB & \cakelpr{}, 4--16 GB heap \\
Historical & 96 & archived DRAT, converted to LRAT by drat-trim; 0.10 TB, largest 2.49 GB & \cakelpr{} \\
Hardest orbit & 1 & 36 cube proofs, 47.1 GB in all & \cakelpr{}; composed by a Lean lemma \\
\midrule
Total & 13,351 & about 28 TB & \\
\bottomrule
\end{tabularx}
\end{table}

\paragraph{The bank.} For 12,094 orbits, the first production run in August 2026 archived a compressed LRAT proof, obtained from a Kissat DRAT proof by drat-trim and then compacted. These proofs had been checked only by drat-trim when produced. We re-checked all of them: regenerate the CNF as above and require its hash to match the producer's ledger; stream the archived file through sha256, decompression, a second sha256 and \cakelpr{}; and accept only if \cakelpr{} verifies and both stream hashes match the ledger. All 12,094 verified, with 22.55 TB of LRAT (6.29 TB compressed), the largest proof 26.5 GB, every check within an 8 GB heap, in 408 hours of summed check time. One orbit has an archived proof but no producer ledger line; it is verified from the pinned-emitter CNF and \cakelpr{} alone, since the ledger match is provenance and not part of the correctness argument. Twenty-one runs reported a CNF mismatch; all were emitter failures on one cloud machine while it was being reclaimed, and every affected orbit re-ran to a verified receipt.

\paragraph{The fresh census.} Of the 1,161 orbits with no August certificate, the 1,160 other than the hardest were re-solved with CaDiCaL 3.0.1 (flags \texttt{-{}-lrat=true -{}-binary=true}) and checked as they streamed by \cakelpr{}: 1,157 on AWS Graviton machines with the published Linux arm64 solver binary (sha256 \texttt{fd601b82\ldots}) and a \cakelpr{} build with sha256 \texttt{95b64883\ldots}, and 3 on a local Apple-silicon machine. All 1,160 verified: 5.64 TB of proofs, 3,435 solver and 207 checker CPU-hours. Of 1,196 run records, 1,163 are verified receipts (three orbits were run twice) and 33 are infrastructure faults (a race in a process watchdog, a 24-hour solver time cap on the longest orbit, and a duplicated claim), each re-run to a verified receipt.

\paragraph{The historical certificates.} The 96 orbits that the August run had also certified were re-checked from their archived DRAT proofs: drat-trim converted each to LRAT against the CNF with the expected hash, and \cakelpr{} (a macOS build of the same commit and assembly file, sha256 \texttt{d23c413b\ldots}) accepted all 96, 0.10 TB of LRAT, the largest proof 2.49 GB. Every converted file has the same sha256 as in an earlier conversion pass. The census and the bank re-check cost about \$225 of cloud compute in total.

\subsection{The hardest orbit}\label{sec:cube}

One orbit, \lean{h1_81494a6ef36d3ec9}, defeated every whole-instance solve, at caps of up to 24 hours. An adaptive cube-and-conquer procedure split it, by unit-propagation lookahead, into a tree of 36 leaf cubes of depth at most 8. Each split replaces a cube by two that partition it, so the leaves partition the formula.

For the certificate, the 36 leaf formulas were re-solved with binary LRAT logging by the macOS build of CaDiCaL 3.0.1, each leaf's CNF hash-checked against the tree. The proofs were written to files and checked afterwards by that macOS \cakelpr{} build, with a 4 GB heap; all 36 verified (47.1 GB in all, the largest 8.24 GB). The leaf receipts record each proof's length, but a proof sha256 only for the 30 smaller leaves (up to 2.42 GB).

The Lean module \repofile{proofs/Proofs/Erdos85H1CubePilot81494a.lean}{Erdos85H1CubePilot81494a.lean} states the tree and the 36 cubes as Lean terms, proves by kernel \texttt{decide} that the cubes cover every cube of the tree (\lean{h1Cube81494a_cover}), and proves \lean{cnf_unsat_of_h1Cube81494a}: the orbit's Lean formula is unsatisfiable if all 36 cube formulas are. The generic step \lean{cnf_unsat_of_cubeTree} and its H1 form \lean{oneHighFamilyV2CheckedUnsat_of_cubeTree} are in \repofile{proofs/Proofs/Erdos85CubeTreeComposition.lean}{Erdos85CubeTreeComposition.lean}. All four statements depend only on \texttt{propext} and \texttt{Quot.sound} (\repofile{\rp/h1_cert_pilot_20261001/COMPOSITION_AXIOMS_20261006.txt}{COMPOSITION_AXIOMS_20261006.txt}). The leaf hypotheses are supplied by the external checks, so for this orbit the composition is checked in Lean and the leaves outside it.

\subsection{What the check establishes}\label{sec:trust}

For each of the 13,351 orbits, \cakelpr{} has accepted a refutation of a CNF file (for the hardest orbit, of 36 cube files) whose hash matches the output of the compiled Lean emitter on that orbit's table. The solvers are no longer trusted. The checks run outside Lean: \cakelpr{}'s guarantee, which covers parsing the files, is a HOL4 theorem about its compiled binary, not a Lean kernel check, and the checks are not admitted into Lean as \lean{OneHighFamilyV2CheckedUnsat} facts. What remains trusted is \cakelpr{}'s HOL4 specification and compiled binary, and the identity between each checked file and the Lean term \lean{oneHighFamilyV2SatCnf}, which rests on the compiled emitter and its printing rather than on a Lean proof. A certificate shows that the checked formula is unsatisfiable, and nothing more. Whether the formula encodes the stratum correctly is a separate question; in H1 the Lean reduction of Section~\ref{sec:h1-structure} answers it for the Lean term, which is why the emitter is the Lean function itself. A reader can re-check any H1 orbit on their own cloud machine with a public machine image, the published bank proofs and a checker kit (Section~\ref{sec:availability}).

\section{Status}\label{sec:status}

Table~\ref{tab:status} is the one place where the paper lists what is checked and what remains open in Lean; the rest of the paper refers to it.

\begin{table}[htbp]
\centering\small
\caption{Status of each component of Result A and of Theorem B. ``Lean'' means a statement elaborated by Lean 4.31.0 from source; where axioms are listed, they are those printed by \texttt{\#print axioms}.}
\label{tab:status}
\begin{tabularx}{\textwidth}{@{}>{\raggedright\arraybackslash}p{.17\textwidth}>{\raggedright\arraybackslash}X>{\raggedright\arraybackslash}p{.31\textwidth}@{}}
\toprule
Component & How it is checked & Open in Lean \\
\midrule
Lower sides & Lean, explicit graphs via \texttt{native\_decide}; cold audit lists three \texttt{native\_decide} axioms per witness & nothing; not standard-axiom-only \\
Case split, $h=9$ & Lean (\leantab{not_c4FreeMinDegreeWitness_fortyNine_seven_of_strata}) & nothing \\
H1 reduction & Lean: 13,351 orbit formulas suffice (\leantab{orderFortyNineStratumExcluded_one_of_capacityInventory_checked}); printed axioms: the three standard ones and 23 \texttt{native\_decide} axioms, no \texttt{sorryAx} & nothing; not standard-axiom-only \\
H1 orbits & 13,351 refutations accepted by \cakelpr{} outside Lean (the hardest as 36 cube leaves, composed by a standard-axiom Lean lemma) & (i) the external checks are not admitted into Lean; (ii) that each checked file is the Lean formula rests on the compiled emitter \vtwocnf, not on a kernel proof \\
H3 (2 cells) & Lean: every graph is covered by the two cell representatives; their exclusion by a reviewed argument with independently replayed computation; no certificates & two cell-exclusion premises \\
H5 (3 cells) & reviewed reductions with independently checked computation; no certificates & three Boolean-exclusion premises \\
H7, $t\ge 1$ (13 cells) & LRAT certificates checked inside Lean by \texttt{native\_decide} & depend on \texttt{Lean.ofReduceBool}; not rebuilt in the cold audit; certificates unpublished \\
H7, $t=0$ & reviewed closure ledger with independently checked computation & capstone not instantiated \\
Theorem B & Lean, standard axioms only & $\AREG$ is a hypothesis \\
\bottomrule
\end{tabularx}
\end{table}

The right-hand column is the complete list of what separates Result A from a Lean theorem, together with a cold build and axiom audit of the assembled statement, which would still list the \texttt{native\_decide} axioms of the witnesses, of the H7 certificates and of the H1 enumeration. For H1 both items are needed: admitting the checks alone would leave open whether the checked files are the Lean formulas. The H3 and H5 premises and the H7 $t=0$ capstone are substantial formalization work; H3 could instead be closed by a certificate for each of its two cell formulas, but neither has been produced, and whether such certificates are within reach is untested.

\section{Theorem B and the residue beneath A-REG}\label{sec:theoremB}

\paragraph{The reduction.} A pair of adjacent orders gives a drop when it has an existence half and a nonexistence half: a $q$-regular $C_4$-free graph on $q^2-1$ vertices, together with the nonexistence of a $C_4$-free graph of minimum degree $q$ on $q^2$ vertices, gives $f(q^2)<f(q^2-1)$ (\lean{PlaneOrderDropWitness.strict_drop}), and cofinally many such pairs refute the problem (\lean{not_erdos85Question_of_cofinalPlaneOrderDropFamily}). For $q=2^k$ with $k\ge 3$ the existence half is uniform: deleting the absolute nucleus from the even-characteristic polarity graph gives the witness (\lean{Polarity.c4FreeMinDegreeWitness_even_delete_absolute_nucleus}). On the nonexistence half the normalized square-order reduction is an equivalence (\lean{squareOrderTightCoreExists_iff_witness}), parity forces every tight core to be regular (\lean{squareOrder_regular_of_even}), and Lean proves that $\AREG$ implies the tight-core exclusion (\lean{binarySquareOrderTightCoreExclusion_of_regularExclusion}) and hence $\neg\,\Qn$. Every arrow is a dependency of Theorem B and was compiled in the cold audit. Theorem B is a complete checked reduction and not a proof of $\AREG$; $\AREG$ does not hide further assumptions.

\paragraph{The residue.} For a hypothetical $q$-regular $C_4$-free graph on $q^2$ vertices with adjacency matrix $A$, the \emph{defect} $D$ is defined by $A^2=(q-1)I+J-D$, and $A$ commutes with $D$. The components of the defect graph have orders $qm_c$ with $\sum_c m_c=q$, and every vertex of component $c$ has exactly $m_c$ neighbours inside it. For even $q$ there are no components with $m_c=1$, and for $4\mid q$ no bipartite components. The Lean statements are \lean{adjMatrix_sq_eq_sub_secondOrderDefect_of_regular}, \lean{binarySquare_regular_exists_defectComponent_partition}, \lean{binarySquare_regular_no_sizeQ_defectComponent_of_even} and \lean{binarySquare_twoPow_regular_defectGraph_not_bipartite}. What remains is $\NONBIP$: partitions $q=\sum_c m_c$ with every $m_c\ge 2$ and every defect component non-bipartite. Its one-component case has a sharp form.

\begin{nonbipconn}
For binary $q\ge 8$, every loopless $q$-regular $C_4$-free adjacency matrix of order $q^2$ is singular. Equivalently, its defect graph is disconnected, since $\dim\ker A=\#\mathrm{components}(D)-1$ (\lean{binarySquare_regular_finrank_adj_kernel_eq_card_components_sub_one}).
\end{nonbipconn}

\noindent The Lean target is \lean{BinarySquareConnectedNonbipartiteExclusion}. There is no uniform Lean reduction of the disconnected cases to the connected one, so the connected case and the mixed non-bipartite partitions are separate open cases. Singularity cannot follow from regularity or the spectrum alone; the frontier is to derive it from the full symmetric $C_4$-free completion.

\paragraph{Evidence on A-REG.} The evidence is mixed. In its favour, the existence half is uniform and the defect calculus removes unit and bipartite components. Against it, Lean proves $f(15)=f(16)=5$ (\lean{minDegreeForC4_fifteen}, \lean{minDegreeForC4_sixteen}), so there is no drop from 15 to 16, and the explicit graph \lean{sixteenRegular} in \repofile{proofs/Proofs/Erdos85Problem.lean}{Erdos85Problem.lean} is 4-regular (\lean{sixteenRegular_degree}) and $C_4$-free (\lean{sixteenRegular_common_le_one}) on 16 vertices: the $q=4$ analogue of $\AREG$ is false. Every generic attempt at the remaining non-bipartite case has failed (Section~\ref{sec:negative}). The principal rival reading is that special plane orders organize both constructions and obstructions, with no uniform binary exclusion. We hold $\AREG$ as a hypothesis whose truth the finite data do not decide.

\section{What does not close A-REG}\label{sec:negative}

Several natural approaches to $\NONBIP$ fail for identifiable reasons, and the failures say what a proof must use: the simultaneous completion of all rows to one symmetric $q$-regular $C_4$-free matrix. The full record is in the campaign's proof outline (\repofile{\rp/FINAL_PROOF_OUTLINE.md}{FINAL_PROOF_OUTLINE.md}) and cuts ledger (\repofile{\rp/CUTS_LEDGER_DRAFT.md}{CUTS_LEDGER_DRAFT.md}); the arguments are reviewed prose unless a Lean name is given.
\begin{itemize}
\item \emph{Determinant and spectrum.} The identity $\det(A)^2=q^4\tau(D)$, with $\tau$ the spanning-tree count, makes a square tree count necessary when $D$ is connected, but controls satisfy it without being singular; real spectral conditions likewise admit controls that are not realizable by integral matrices.
\item \emph{Inverse signs.} A proposed separation of the signs of inverse entries (P8) is impossible: exact product-sign identities force both signs. It survived a thousand sampled $q=4$ models only because those models were singular and never instantiated its hypothesis.
\item \emph{Structured families.} The interval circulant defect family is excluded for every binary $q\ge 4$ by a cyclotomic argument (\repofile{\rp/l6-cyclic-trace/UNIFORM.md}{l6-cyclic-trace/UNIFORM.md}), and specific order-256 defects fail local Hasse invariants; neither reduces arbitrary defect graphs to those families. Explicit connected non-bipartite defect graphs pass the alternating square-root test over $\mathbb{F}_2$ but fail the determinant-square condition (\repofile{\rp/mod2-root-gate/}{mod2-root-gate/}).
\item \emph{Relaxations.} Packing, Hall and fractional-transversal relaxations omit the global completion, and the strongest faithful partial version already has an exact half-integral survivor at $q=6$; a five-vertex flag relaxation is feasible for every binary $q\ge 16$ (cuts ledger, row 175).
\item \emph{Order 64.} At $q=8$ some defect partitions are closed in Lean, including the size-two sector (\lean{orderSixtyFour_regular_sizeTwo_signedJoint_false_of_connected}), but $[3,3,2]$, $[4,2,2]$ and $[6,2]$ keep non-bipartite cases, $[4,4]$ and $[5,3]$ are only partly analysed and the connected $[8]$ case is open, so 63 to 64 is not a drop.
\item \emph{Plane orders 9, 11, 13.} The order-80 search for a 9-regular $C_4$-free graph ended UNKNOWN in all 20 solver attempts (\repofile{\rp/q9_solver_controls/Q9_EXISTENCE_DECISION_20260911.md}{Q9_EXISTENCE_DECISION_20260911.md}). A census of Cayley graphs on 52 groups of order 80, 47 of order 120 and 57 of order 168 found no $C_4$-free Cayley graph of degree 9, 11 or 13 (\repofile{\rp/cayley-census-q11-q13/CAYLEY_CENSUS_Q11_Q13_20260913.md}{CAYLEY_CENSUS_Q11_Q13_20260913.md}); this excludes only those families. With the values of \citet{zhang2017polarity} and \citet{boza2024ramsey} it leaves $r(109)\in\{120,121\}$ and $r(155)\in\{168,169\}$; our literature check found no exact value for either within its stated scope.
\end{itemize}

\section{Collaboration and methods}\label{sec:collab}

\paragraph{Contributions.} Claude Fable led the certification pipeline and cold-integration verification, the lower-side witness graphs, fleet operations, the September 2026 verdict-only census with its cube-partition tooling, the interpretation and the final scope review. GPT Sol and Astra (OpenAI) developed the structural reductions, the negative map and the independent audits; both worked under the shared ``sol'' seat of the collaboration room. Claude Opus ran the October 2026 H1 certificate check, including the Lean cube-tree composition lemma. Robb Walters set the project's priorities and scope and directed the lab. The shared room infrastructure, which supplied coordination, claims, reviews and a durable transcript (\url{https://rjwalters.info/rooms/erdos-85}), is acknowledged; it is not an author.

\paragraph{Methods.} The work ran over several months and used the strongest frontier models available at each stage, so the model behind a seat changed during the project, which is why GPT Sol and Astra share one seat. Proposals were cheap; admission to the record was not. A Lean result counted only after it elaborated from source under the pinned toolchain in an independent environment and its \texttt{\#print axioms} output was read, a rule adopted after a stale-cache incident. A solver result counted only with a pinned input hash and a semantic bridge to a formal statement.

\paragraph{Two lessons.} \emph{The encoding is the weak link.} At order 64, a first encoding of one endpoint silently reused the shore table of a neighbouring case (80 candidate owners instead of 88) and returned UNSAT for a strengthened, wrong problem. The mismatch was found by comparing the formal shore modes with the certificate universe before the endpoint was called proved; the corrected 88-owner formulas are checked in \repofile{proofs/Proofs/Erdos85MuNegThreeZeroFiveCorrectOwnerCertificate.lean}{Erdos85MuNegThreeZeroFiveCorrectOwnerCertificate.lean}. The H1 pipeline guards against the same failure by emitting each formula with the Lean function that the Lean reduction quantifies over. \emph{Silence is not success.} \cakelpr{} reports a failed check with exit status 0, the P8 separation survived only on models that never met its hypothesis, and long builds that end without an error have not yet printed their axioms. In each case the remedy was to require a positive artifact: a verification line, an instantiated hypothesis, a printed axiom list.

\section{Availability}\label{sec:availability}

Everything named here is in the Lean Genius repository (\url{https://github.com/rjwalters/lean-genius}) on the branch \texttt{erdos85/integration}, under \repofile{\rp/}{research/problems/erdos-85-wip-01/} and \repofile{proofs/Proofs/}{proofs/Proofs/}; labels below are relative to those directories.
\begin{itemize}\raggedright
\item H1 bank re-check: \repofile{\rp/h1_bank_check_20261006/receipts/README.md}{h1_bank_check_20261006/receipts/README.md}, with one receipt per orbit in \repofile{\rp/h1_bank_check_20261006/receipts/h1_bank_check_receipts.tsv}{h1_bank_check_receipts.tsv} and totals in \repofile{\rp/h1_bank_check_20261006/receipts/h1_bank_check_summary.json}{h1_bank_check_summary.json}; tooling in \repofile{\rp/h1_bank_check_20261006/}{h1_bank_check_20261006/}.
\item Fresh census, historical certificates and cube leaves: \repofile{\rp/h1_cert_full_20261001/receipts/README.md}{h1_cert_full_20261001/receipts/README.md} (pinned tool hashes and a reproduction recipe), \repofile{\rp/h1_cert_full_20261001/receipts/h1_cert_census_receipts.tsv}{h1_cert_census_receipts.tsv}, \repofile{\rp/h1_cert_full_20261001/receipts/h1_cert_census_summary.json}{h1_cert_census_summary.json}, \repofile{\rp/h1_cert_full_20261001/receipts/historical96_cake_lpr_receipts.jsonl}{historical96_cake_lpr_receipts.jsonl}, \repofile{\rp/h1_cert_full_20261001/receipts/historical96_receipts.jsonl}{historical96_receipts.jsonl} (the earlier conversion pass), \repofile{\rp/h1_cert_full_20261001/receipts/pilot_h1_81494a_leaf_cake_lpr.jsonl}{pilot_h1_81494a_leaf_cake_lpr.jsonl}; tooling in \repofile{\rp/h1_cert_full_20261001/}{h1_cert_full_20261001/} and \repofile{\rp/h1_cert_pilot_20261001/}{h1_cert_pilot_20261001/}; axiom lists of the H1 cover theorem and of the cube composition in \repofile{\rp/h1_cert_pilot_20261001/H1_COVER_AXIOMS_20261006.txt}{H1_COVER_AXIOMS_20261006.txt} and \repofile{\rp/h1_cert_pilot_20261001/COMPOSITION_AXIOMS_20261006.txt}{COMPOSITION_AXIOMS_20261006.txt}.
\item Verdict census and cube tree: \repofile{\rp/phase_b_h1_census_20260927/CENSUS.md}{CENSUS.md}, \repofile{\rp/phase_b_h1_census_20260927/cube-h1_81494a6ef36d3ec9/tree-check.json}{tree-check.json}; core-hours in \repofile{\rp/CENSUS_TIMING_20260928.md}{CENSUS_TIMING_20260928.md}; the H1 inventory and the 12,094 bank orbits in \repofile{\rp/PHASE_B_H1_H3_INVENTORY_20260910.md}{PHASE_B_H1_H3_INVENTORY_20260910.md}.
\item Lean: the H1 reduction in \repofile{proofs/Proofs/Erdos85OneHighV2CapacityCover.lean}{Erdos85OneHighV2CapacityCover.lean} and \repofile{proofs/Proofs/Erdos85OneHighV2CapacityInventory.lean}{Erdos85OneHighV2CapacityInventory.lean}, the inventory file \repofile{proofs/Proofs/Certificates/h1_orbit_inventory.compact}{Certificates/h1_orbit_inventory.compact}; the cold axiom audit in \repofile{\rp/AXIOM_AUDIT_COLD_20260927/axioms.out}{AXIOM_AUDIT_COLD_20260927/axioms.out}; the strata and small-order names in \repofile{\rp/STRATA_AND_SMALL_ORDERS_20260928.md}{STRATA_AND_SMALL_ORDERS_20260928.md}.
\item Small strata: \repofile{\rp/Q7_H3_PROFILE_EXCLUSION_20260910.md}{Q7_H3_PROFILE_EXCLUSION_20260910.md} (H3), \repofile{\rp/H7_POSITIVE_TRIPLE_CELLS_H5_CLOSURE_FORMULA_COUNTS_20260928.md}{H7_POSITIVE_TRIPLE_CELLS_H5_CLOSURE_FORMULA_COUNTS_20260928.md}, \repofile{\rp/q7_h5_closure_ledger/REVIEWED_RESULT.md}{q7_h5_closure_ledger/REVIEWED_RESULT.md}, \repofile{\rp/H7_CLOSURE_20260915.md}{H7_CLOSURE_20260915.md} with \repofile{\rp/h7-closure-ledger-20260915/}{h7-closure-ledger-20260915/}, and \repofile{\rp/Q7_H5_H7_SQUEEZE_20260910.md}{Q7_H5_H7_SQUEEZE_20260910.md}.
\item Literature readings and the $r$-to-$f$ conversion: \repofile{\rp/manuscript/FIRST_DROP_LITERATURE_CHECK.md}{manuscript/FIRST_DROP_LITERATURE_CHECK.md}; the order-48 comparison receipt \repofile{\rp/sat49/BOZA48_NONISOMORPHISM_RECEIPT_20260928.txt}{sat49/BOZA48_NONISOMORPHISM_RECEIPT_20260928.txt}.
\end{itemize}
\paragraph{Checking H1 yourself.} A reader can re-check any H1 orbit on their own AWS machine without downloading anything locally. The public machine image \texttt{ami-05697724475f2e748} (region us-east-1) contains the pinned emitter, solver and \cakelpr{} and the \texttt{e85-check} tool; the 12,094 bank proofs and the checker kit (manifests, orbit tables, our receipts and pinned binaries) are in an Amazon S3 bucket with Requester Pays access. \texttt{e85-check} streams a bank proof into \cakelpr{}, re-solves a census or historical orbit with the pinned CaDiCaL and checks the fresh proof with \cakelpr{} (for the 1,157 census orbits we solved on Linux the proof's sha256 must also equal ours; for the 96 historical orbits and the 3 census orbits solved on macOS there is no comparable proof hash, so the result is a fresh \cakelpr{} check), or rebuilds and checks the cube leaves of the hardest orbit. The guide \repofile{\rp/CHECKING.md}{CHECKING.md} gives the steps, the bucket paths and the meaning of each result, and \repofile{\rp/h1_checker_kit/}{h1_checker_kit/} holds the tooling.

The fresh H1 proofs were checked and discarded; they are regenerated by running the pinned CaDiCaL binary on the pinned CNF and compared by sha256. Not published: the rest of the certificate storage, the packed H7 certificates, and the artifact volume with full solver logs and run directories.

\bibliographystyle{plainnat}
\bibliography{refs}

\end{document}
